% ChargeShield-FL --- DSN 2027 --- Architecture / Framework

\section{Architecture}
\label{sec:framework}

\begin{quote}
\emph{A terminological clarification, stated once here because it is
used without re-explanation throughout the rest of the paper.} We use
\emph{auditing} for two distinct activities and keep them separate by
name. \emph{Budget auditing} is the in-line activity of the Privacy
Auditor (Section~\ref{sec:framework-auditor}): it accounts for what the
DP mechanism \emph{promises}. \emph{Empirical privacy auditing} is the
offline activity of the attack suite (Section~\ref{sec:framework-attacks}):
it measures what an instantiated adversary actually \emph{extracts}.
These are different components, running at different times, serving
different roles; neither substitutes for the other.
\end{quote}

\subsection{The ML Plane}
\label{sec:framework-mlplane}

The ML Plane is an event-driven observability substrate for FL
training. Each component that participates in a training round --- the
local trainer, the gradient/DP handler, the FedAvg aggregator ---
exposes a common \texttt{emit\_event()}/\texttt{subscribe()} interface;
a central hub connects publishers to subscribers; a collector
assembles, for every round, both the raw and the privatized update,
without the training loop itself knowing that anything is listening.
Algorithm~\ref{alg:mlplane-dispatch} gives the dispatch loop.

Three event types exist, and they are also the three observation
points used in the experiments of Section~\ref{sec:results}: \emph{local
update produced}, \emph{privatized update}, \emph{aggregated round}.
Every event carries the Purdue level at which the artifact was
produced --- level 1 for the raw local-training update, level 2 for
the clipped-and-noised update that reaches the supervisory level. This
label is not decorative: it determines which artifact a given consumer
is entitled to see, and it is what makes expressible the difference
between a deployment where the aggregator observes the clean value and
one where it never does.

Nothing at this layer is specific to the EV domain or to any
particular attack. Figure~\ref{fig:mlplane} shows the three consumers
introduced in Sections~\ref{sec:framework-auditor}--\ref{sec:framework-attacks}
attached to the same event stream.

\begin{algorithm}[t]
\caption{ML Plane event dispatch (per training round)}
\label{alg:mlplane-dispatch}
\begin{algorithmic}[1]
\REQUIRE round index $t$, set of participating clients $C$
\STATE $U_{\text{raw}} \gets \{\}$, $U_{\text{priv}} \gets \{\}$
\FOR{each client $c \in C$}
  \STATE $u_c \gets \textsc{LocalTrain}(c, t)$ \textit{Purdue level 1}
  \STATE \textsc{Emit}(\textsc{LocalUpdateProduced}, $u_c$, level=1)
  \STATE $\hat{u}_c \gets \textsc{ClipAndPrivatize}(u_c)$ \textit{Purdue level 2}
  \STATE \textsc{Emit}(\textsc{PrivatizedUpdateReady}, $\hat{u}_c$, level=2)
  \STATE $U_{\text{raw}} \gets U_{\text{raw}} \cup \{u_c\}$; \quad $U_{\text{priv}} \gets U_{\text{priv}} \cup \{\hat{u}_c\}$
\ENDFOR
\STATE $g_t \gets \textsc{FedAvgAggregate}(U_{\text{priv}})$
\STATE \textsc{Emit}(\textsc{RoundAggregated}, $g_t$, $U_{\text{raw}}$, $U_{\text{priv}}$)
\STATE \textit{Every subscriber below receives every emitted event; none blocks or mutates the training loop.}
\STATE \textbf{subscribers:} PrivacyAuditor.\textsc{OnEvent}($\cdot$) \COMMENT{\S\ref{sec:framework-auditor}, in-line}
\STATE \hphantom{\textbf{subscribers:}} ByzantineDetector.\textsc{OnEvent}{$\cdot$} \COMMENT{\S\ref{sec:framework-byzantine}, in-line}
\STATE \hphantom{\textbf{subscribers:}} RoundDataCollector.\textsc{OnEvent}{$\cdot$} \COMMENT{\S\ref{sec:framework-attacks}, buffers only; attacks run later, offline}
\end{algorithmic}
\end{algorithm}

\subsection{First Consumer --- the Privacy Auditor (in-line telemetry)}
\label{sec:framework-auditor}

The Privacy Auditor is an always-on, in-line consumer subscribed to ML
Plane events. Given any node's update as a flat dictionary of numeric
weights, it does exactly three things: it computes a sensitivity proxy
based on the update's norm, it tracks cumulative $\varepsilon/\delta$
budget consumption per node across rounds, and it raises alerts when
configured thresholds are crossed (gradient explosion, budget near
exhaustion or exhausted).

\textbf{It has no notion of membership and no notion of attack.} It
never evaluates whether a specific session belongs to a node's training
set, never computes a membership score, and never runs any of the
attacks in Section~\ref{sec:framework-attacks}. It is privacy
\emph{telemetry}, not a privacy \emph{measurement}: it answers ``is
this node behaving within its declared parameters?'', not ``how much
information is actually leaving?''

We state two qualifications explicitly.

\emph{The budget score is not formal DP accounting.} Per-round
$\varepsilon$ is computed as
$(\text{observed sensitivity}/\texttt{max\_grad\_norm}) \times
\varepsilon_{\text{budget}}$, where observed sensitivity is scaled
relative to the round's peer median (Algorithm~\ref{alg:round-epsilon}).
It is therefore \textbf{data-dependent}, whereas formal DP accounting
is data-independent by construction, and it can exceed the nominal
per-round budget when a client sits above the median: in our exported
artifacts one client reports $\varepsilon = 1.070$ against a budget of
$1.0$ at the first round. This is an internal risk score calibrated
for alert thresholds, not an accountant, and we disclose it both here
and in the \texttt{epsilon\_accounting\_note} field of every exported
JSON report, because a reader who opens only the artifact cannot
otherwise infer it.

\emph{The Auditor has never been validated against a true positive.}
Across the full experimental campaign it has raised no alert, because
there has never been anything to flag. A monitor that has never seen
the event it is meant to detect has not been evaluated as a monitor.
The most direct way to close this gap is to point the Auditor at a run
where membership leakage is known and deliberately induced --- i.e.,
at the canary positive-control runs described in
Section~\ref{sec:framework-canary} --- and check whether its
sensitivity proxy moves. We flag one dependency before that check can
be treated as conclusive: a tagging bug affecting the non-member side
of canary injection (an untagged original session could land inside a
shadow model's training set while its tagged clone was held out, or
vice versa, defeating the anti-contamination guard the tagging exists
to enforce) was found and fixed on 2026-09-15, after every canary run
reported in an earlier draft of this paper had already completed;
those runs should be read as indicative of the Auditor-validation
methodology, not as its result, pending a rerun against the fixed
injection code. Overhead --- wall-clock time spent per round inside the
Auditor's own processing, separate from training time --- is
instrumented (a single always-on timer per round, chosen over a
two-run A/B comparison because cross-run training-time variance would
dominate an expected sub-millisecond signal) but not yet reported here
pending a run under realistic DP configuration.

\subsection{Second Consumer --- ByzantineDetector (in-line integrity)}
\label{sec:framework-byzantine}

ByzantineDetector (Krum~\cite{blanchard2017machine},
cosine-similarity filtering, CUSUM, gradient-explosion detection)
protects training convergence against Byzantine or poisoning clients.
It is the ML Plane's second in-line consumer, and it operates on a
threat domain orthogonal to the Auditor's: client-side integrity
(Scenario~2, Section~\ref{sec:threat-model}) versus server-side
confidentiality (Scenario~1, Section~\ref{sec:threat-model}).

We state explicitly, having found and corrected an earlier conflation
of our own on this point: \textbf{no integrity mechanism provides any
privacy guarantee} against the honest-but-curious aggregator this
paper's results characterize, and no privacy measure detects Byzantine
behavior. The two domains are separated in the architecture and remain
separated throughout our argument.

As of this submission, ByzantineDetector's mechanisms are implemented
and unit-tested, but have not been exercised end-to-end against an
active Byzantine client on real training data.

\emph{An observation worth reporting if that experiment is run.} A
preliminary numerical simulation, independent of the production
ByzantineDetector code, suggests that DP noise saturates the detector:
with $\sigma = 4.84$ (corresponding to $\varepsilon = 1.0$ under our
noise calibration, Section~\ref{sec:framework-dp}) applied to two
simulated \emph{honest} updates of dimension matching our
570-parameter model ($w \sim \mathcal{N}(0, 0.1)$ per weight before
noise), mean cosine similarity between the two updates collapses from
$\approx 0.9999$ in a noise-free control to $\approx -0.0024$ after
noise, with $100\%$ of 200 simulated honest pairs falling below
ByzantineDetector's default $0.3$ similarity threshold. If confirmed on
real trained weights, this is a concrete privacy/integrity tension:
noise calibrated for confidentiality renders integrity monitoring on
the same channel inoperative, independent of DP placement
(Section~\ref{sec:framework-dp}). This would be a directly actionable
finding for deployment design.

\subsection{Third Consumer --- the Attack Suite (offline, empirical auditing)}
\label{sec:framework-attacks}

A $(\varepsilon,\delta)$ guarantee bounds worst-case leakage over every
possible adversary and every neighboring dataset; it does not say what
a concrete adversary extracts from one specific, already-trained
model. Measuring the latter quantity requires instantiating the
adversary. That is the role of the attack suite.

The attacks --- Yeom, Shadow, LiRA, in \texttt{src/plugins/attacks/},
dispatched through \texttt{ATTACK\_REGISTRY} --- consume the same
\texttt{round\_data} collected by the ML Plane, but run
\textbf{offline, after training completes}, as a separate process. In
the single-process simulation, \texttt{run\_experiments.py} finishes
training and then iterates the registry. In the NVFlare deployment,
\texttt{ChargeShieldAggregator} exports a round-data dump at every
round and \texttt{run\_nvflare\_mia.py} loads it and invokes the same,
unmodified attack functions. In neither path does attack code run
inside the server process. Section~\ref{sec:framework-offline-equiv}
justifies this as a design choice on equivalence grounds, not
implementation convenience.

A fourth, opt-in diagnostic attack, FedMIA-gradient, reuses two
numerical routines (\texttt{calibrate\_from\_vectors()},
\texttt{reconstruction\_error()}) from a historical, never-activated
\texttt{FedMIA} class in this codebase, but implements a distinct
gradient-space reconstruction-error test; it is not validated to the
same standard as the other three attacks and is reported only as a
diagnostic (Section~\ref{sec:limitations}). See
Section~\ref{sec:related-mia} for why neither artifact should be read
as an implementation of Zhu~et~al.'s unrelated, identically-named
method~\cite{zhu2025fedmia}.

\subsubsection{Canary Positive Control}
\label{sec:framework-canary}

To validate that a null attack result reflects the absence of leakage
rather than an insensitive instrument, we inject synthetic canary
sessions --- near-duplicate templates repeated with high multiplicity
--- into a client's training set and verify that the attack suite
\emph{can} detect them when a true membership signal is deliberately
introduced. Algorithm~\ref{alg:inject-canaries} gives the injection
procedure, including the tagging discipline required on both the
member and non-member side: each canary template's original,
untagged occurrence must be replaced in place by a tagged copy before
any duplicate is added, on \emph{both} sides of the split, so that
downstream shadow-model training (which must never train on a sample
whose exact duplicate it is later asked to score) can identify and
exclude every copy of a template consistently. An asymmetric version of
this bug --- present on the member side only, then independently found
and fixed on the non-member side as well on 2026-09-15 --- is the
subject of the caveat in Section~\ref{sec:framework-auditor}.

Two properties of this design bound what the control can establish, and
we state them explicitly because the nominal sample size overstates
both. First, \textbf{effective sample size}. With $|T_m|$ member
templates each repeated $M+1$ times and $|T_n|$ non-member templates
each present twice, every copy of a template carries an identical
feature vector, hence an identical reconstruction loss and --- because
the group is held atomic across shadow splits, as required to prevent
the contamination described above --- an identical calibrated score.
The AUC computed over the $(M{+}1)|T_m| \times 2|T_n|$ nominal pairs is
therefore numerically \emph{equal} to the AUC over the
$|T_m| \times |T_n|$ \emph{distinct} pairs, but its sampling variance
is that of the smaller quantity. In our configuration
($|T_m| = 5$, $M = 30$, $|T_n| = 20$) a nominal $n = 155 \times 36$
corresponds to at most $5 \times 20$ independent comparisons: a
confidence interval or significance test computed on the nominal $n$ is
narrower than the data support by a factor of $\sqrt{62} \approx 7.9$.
We therefore read canary AUC as a \emph{direction} indicator and never
as a precise effect size, and we attach no $p$-values to it. The
signature is visible in the raw outputs: every canary AUC we obtain is
an exact multiple of $1/(|T_m|\,|T_n^{\text{eff}}|)$, never of the
nominal reciprocal.

Second, \textbf{a role-exchange control requires balanced groups}.
Re-running one seed with the member and non-member roles exchanged is
the decisive check that a canary AUC above $0.5$ reflects membership
rather than the intrinsic reconstruction difficulty of whichever
templates the draw happened to select. That exchange is well defined
only when $|T_m| = |T_n|$. With $|T_m| \neq |T_n|$ the two arms draw
\emph{disjoint} member groups from the same pool while sharing most of
the non-member pool --- a second experiment on different records, not a
role exchange --- and its outcome carries no negative-control force. A
companion check measuring the same AUC at random initialisation, before
any training, isolates the intrinsic-difficulty component; it must
reproduce the training run's split procedure exactly, including the
shuffle that precedes the train/hold-out cut, or it characterises a
different set of templates than the one under test. We report both
controls, and the conditions under which each is valid, rather than
reporting the canary AUC alone.

\subsection{Three DP Placements as Adversary Privilege Points}
\label{sec:framework-dp}

The three DP placements we evaluate --- \texttt{dp-fedavg},
\texttt{central}, \texttt{local} --- are not three privacy mechanisms.
They are three points of adversary privilege along the same update
transformation pipeline (Figure~\ref{fig:dp-placements}):
\texttt{dp-fedavg} observes the \emph{raw} update, \texttt{central}
observes it \emph{after clipping but before noise}, \texttt{local}
observes only the \emph{final, clipped-and-noised} value.
\texttt{dp-fedavg} is therefore the \textbf{upper bound} on what is
extractable from this channel. If even the most privileged view shows
no signal, a null result under the protective placements cannot be
attributed to noise suppressing a signal that was otherwise present ---
the raw view already did not contain one.

Two qualifications carry over from this placement design. First, our
clipping and noise operate at \textbf{client-level} granularity (one
update per client per round), while the attacks in
Section~\ref{sec:framework-attacks} probe \textbf{record-level}
membership (one session); the two granularities are not the same
privacy unit, and a client-level guarantee does not by itself bound
record-level leakage when a client holds many records, as is the case
here (each site holds thousands of sessions). Second, $\varepsilon$
here is read as an experimental noise parameter in the presence of 50
local epochs per round, not as a claim about the DP guarantee that
would hold under a single global epoch; the two are numerically
different questions and we do not conflate them.

Noise is calibrated per weight of the 570-parameter autoencoder as
$\sigma = C \cdot \sqrt{2 \ln(1.25/\delta)} / \varepsilon$, with
$C = 1.0$ (the clipping norm) and $\delta = 10^{-5}$, giving
$\sigma = 4.84$ at $\varepsilon = 1.0$ and $\sigma = 48.4$ at
$\varepsilon = 0.1$ (Section~\ref{sec:results-utility}).

\subsection{LiRA: Adaptations to This Setting}
\label{sec:framework-lira}

Our LiRA implementation follows four declared deviations from
Carlini~et~al.'s original construction~\cite{carlini2022membership},
made necessary by the federated, cross-round setting: per-round shadow
retraining from that round's global weights rather than a single
offline shadow population; $n_{\text{shadow}} = 16$ rather than the
hundreds used in the centralized, single-shot setting, a scale
constrained by wall-clock cost per round; three interchangeable
per-sample scorers (raw MSE, log-transformed MSE, and the
Sablayrolles~et~al.\ score~\cite{sablayrolles2019whitebox}); and a
cross-round composed score that accumulates evidence forward rather
than scoring each round independently. We measured, rather than
assumed, the Gaussian-loss assumption underlying the parametric variant:
a Jarque--Bera normality test on raw per-sample loss rejects
Gaussianity at essentially every round (details in
Section~\ref{sec:validation}), motivating the log-transformed variant
as an explicit, declared alternative rather than a silent substitute;
the empirically measured skip rate under the $8\sigma$ exclusion filter
across the full campaign is $0.0000\%$.

\subsubsection{The Shared Variance Floor}
\label{sec:framework-variance-floor}

One deviation is substantial enough to state as its own result rather
than a footnote. When the in- and out-distribution variances are
equal, $\sigma_{\text{in}} = \sigma_{\text{out}} = \sigma$, the two
$-\log\sigma$ normalization terms in the Gaussian log-likelihood ratio
cancel, and the score collapses algebraically to a strictly monotone
\emph{linear} function of the raw loss with a decision threshold at
the midpoint of the two class means:
\begin{equation}
\text{LLR} \;=\; \frac{(\mu_{\text{in}} - \mu_{\text{out}})\,\bigl(2\cdot\text{loss} - \mu_{\text{in}} - \mu_{\text{out}}\bigr)}{2\sigma^2},
\qquad \text{threshold at } \text{loss} = \frac{\mu_{\text{in}}+\mu_{\text{out}}}{2}.
\label{eq:variance-floor}
\end{equation}
Algorithm~\ref{alg:lira-score} shows where this collapse occurs inside
per-round scoring. We verified Equation~\ref{eq:variance-floor} by
independent hand derivation from the Gaussian log-likelihood ratio; it
matches our implementation exactly. In the campaign, the variance floor
is active --- i.e., $\sigma_{\text{in}}$ and $\sigma_{\text{out}}$ are
clamped to the same floor value --- on $95$--$99\%$ of samples under
\texttt{central} and on $68\%$ elsewhere, meaning per-example variance
calibration, the one property that formally distinguishes LiRA from a
Yeom-style loss threshold, is in practice inactive for most of the
campaign. We report this as a declared deviation
(Section~\ref{sec:limitations}) rather than treat LiRA's results as
evidence of the calibrated construction's specific behavior.

\subsubsection{Cold-Start Ablation}
\label{sec:framework-coldstart}

The null result also holds when shadow models are trained from random
initialization (Carlini~et~al.'s original construction) rather than
warm-started from the current global weights, our default. A
collateral observation: under cold-start,
\texttt{matched\_formula\_auc} (Section~\ref{sec:validation}) rises to
$0.49$--$0.51$, compared to $0.34$--$0.44$ under warm-start, while the
variance-floor hit rate rises to nearly $100\%$. These two variables
interact and should be read together, not independently.

\begin{algorithm}[t]
\caption{Canary injection with member/non-member tagging (both sides required)}
\label{alg:inject-canaries}
\begin{algorithmic}[1]
\REQUIRE member templates $T_m$, non-member (holdout) templates $T_n$, duplicate count $d$
\ENSURE tagged injected training set $I_{\text{train}}$, tagged injected holdout set $I_{\text{holdout}}$
\FOR{$j \gets 1$ to $|T_m|$}
  \STATE $\text{template} \gets T_m[j]$; \quad $\text{group} \gets$ \texttt{"canary\_m"} $\Vert\, j$
  \STATE \textbf{find} the occurrence $s$ in $I_{\text{train}}$ with $s \text{ is } \text{template}$ (identity, not equality)
  \STATE replace it in place: $s \gets \text{template} \cup \{\texttt{\_canary\_group}: \text{group},\ \texttt{\_canary\_role}: \texttt{"member"}\}$
  \FOR{$k \gets 1$ to $d$}
    \STATE clone $\gets \text{template} \cup \{\texttt{\_canary\_group}: \text{group},\ \texttt{\_canary\_role}: \texttt{"member"}\}$
    \STATE append clone to $I_{\text{train}}$
  \ENDFOR
\ENDFOR
\STATE \textit{Non-member side --- mirrors the member side exactly. Sprint 10zz+108 (2026-09-15) found this branch previously appended only the tagged clone and left the original \texttt{template} occurrence untagged inside $I_{\text{holdout}}$, reproducing the same shadow-contamination risk this whole procedure exists to prevent: an untagged twin of a template could land IN a shadow model's training set while its tagged clone was OUT, or vice versa.}
\FOR{$j \gets 1$ to $|T_n|$}
  \STATE $\text{template} \gets T_n[j]$; \quad $\text{group} \gets$ \texttt{"canary\_n"} $\Vert\, j$
  \STATE \textbf{find} the occurrence $s$ in $I_{\text{holdout}}$ with $s \text{ is } \text{template}$ (identity, not equality)
  \STATE replace it in place: $s \gets \text{template} \cup \{\texttt{\_canary\_group}: \text{group},\ \texttt{\_canary\_role}: \texttt{"nonmember"}\}$
  \STATE clone $\gets \text{template} \cup \{\texttt{\_canary\_group}: \text{group},\ \texttt{\_canary\_role}: \texttt{"nonmember"}\}$
  \STATE append clone to $I_{\text{holdout}}$
\ENDFOR
\STATE \RETURN $I_{\text{train}}, I_{\text{holdout}}$
\end{algorithmic}
\end{algorithm}

\begin{algorithm}[t]
\caption{LiRA per-round scoring with variance-floor collapse}
\label{alg:lira-score}
\begin{algorithmic}[1]
\REQUIRE target loss $\ell$, shadow-model in/out losses $\{\ell^{\text{in}}_i\}_{i=1}^{n_{\text{shadow}}}$, $\{\ell^{\text{out}}_i\}_{i=1}^{n_{\text{shadow}}}$, variance floor $\sigma^2_{\min}$
\STATE $\mu_{\text{in}}, \sigma^2_{\text{in}} \gets \textsc{Mean}(\ell^{\text{in}}), \textsc{Var}(\ell^{\text{in}})$
\STATE $\mu_{\text{out}}, \sigma^2_{\text{out}} \gets \textsc{Mean}(\ell^{\text{out}}), \textsc{Var}(\ell^{\text{out}})$
\STATE $\sigma^2_{\text{in}} \gets \max(\sigma^2_{\text{in}}, \sigma^2_{\min})$; \quad $\sigma^2_{\text{out}} \gets \max(\sigma^2_{\text{out}}, \sigma^2_{\min})$
\IF{$\sigma^2_{\text{in}} = \sigma^2_{\text{out}}$} 
  \STATE \textbf{score} $\gets \dfrac{(\mu_{\text{in}}-\mu_{\text{out}})(2\ell-\mu_{\text{in}}-\mu_{\text{out}})}{2\sigma_{\text{in}}^2}$ \textit{(linear in $\ell$; threshold at $(\mu_{\text{in}}+\mu_{\text{out}})/2$)}
\ELSE
  \STATE \textbf{score} $\gets \log\dfrac{\mathcal{N}(\ell;\ \mu_{\text{in}}, \sigma^2_{\text{in}})}{\mathcal{N}(\ell;\ \mu_{\text{out}}, \sigma^2_{\text{out}})}$ \COMMENT{full likelihood ratio, per-example variance calibration active}
\ENDIF
\STATE \RETURN \textbf{score}
\end{algorithmic}
\end{algorithm}

\begin{algorithm}[t]
\caption{Privacy Auditor per-round budget score (not formal DP accounting)}
\label{alg:round-epsilon}
\begin{algorithmic}[1]
\REQUIRE client update $u_c$, peer updates $\{u_{c'}\}_{c' \in C}$ for round $t$, \texttt{max\_grad\_norm}, $\varepsilon_{\text{budget}}$
\STATE $m_t \gets \textsc{Median}(\{\lVert u_{c'} \rVert\}_{c' \in C})$ \COMMENT{peer-relative normalization}
\STATE $\hat{s}_c \gets \lVert u_c \rVert \,/\, m_t$ \COMMENT{sensitivity proxy, data-dependent}
\STATE $\varepsilon_c \gets (\hat{s}_c \,/\, \texttt{max\_grad\_norm}) \times \varepsilon_{\text{budget}}$
\IF{$\varepsilon_c > \varepsilon_{\text{budget}}$} 
  \STATE \textsc{RaiseAlert}(\textsc{BudgetExceeded}, $c$, $\varepsilon_c$)
\ENDIF
\STATE \textsc{record} $\varepsilon_c$ in \texttt{epsilon\_accounting\_note}-annotated report for client $c$, round $t$
\STATE \RETURN $\varepsilon_c$
\end{algorithmic}
\end{algorithm}